\qquad Em várias situações do cotidiano é necessário manter conjuntos de
dados organizados: da lista telefônica ao histórico de transações
bancárias, a ordenação é uma operação central em ciência da
computação. Diante disso, diversos algoritmos de ordenação foram
desenvolvidos ao longo dos anos, cada um com características
próprias de desempenho, simplicidade de implementação e
comportamento em diferentes cenários de entrada.

\qquad Este trabalho é dividido em quatro capítulos, um para cada
algoritmo de ordenação estudado: o Capítulo 1 trata do Insertion
Sort, o Capítulo 2 do Selection Sort, o Capítulo 3 do Bubble Sort e
o Capítulo 4 do Shell Sort. Todos são algoritmos de ordenação por
comparação amplamente utilizados como introdução ao estudo de
algoritmos, e, com exceção do Shell Sort, todos possuem complexidade
$O(n^2)$ no pior caso — ainda que, dependendo dos dados de entrada,
alguns deles possam ser bem mais rápidos do que essa cota sugere.

\qquad Ao longo de cada capítulo são apresentados o funcionamento do
respectivo algoritmo, um teste de mesa manual, a análise de
complexidade nas diferentes situações de entrada, e os resultados
experimentais obtidos a partir de uma implementação em C++, testada
com instâncias de tamanhos crescentes (de 10 a 1.000.000 de
elementos) e três padrões de entrada: crescente, decrescente e
randômica. Ao final do trabalho, os quatro algoritmos são comparados
entre si por meio de uma tabela e um gráfico geral, seguidos de uma
conclusão consolidada sobre os resultados observados.
